% !TeX root = closed-systems-from-comparison-completeness.tex
% ============================================================
\chapter{Comparison Completeness and the Emergence of Diagonal Redundancy}
\label{ch:bottom}
% ============================================================
\section{Introduction}
Under the standing principle of closed-world admissibility
(\cref{sp:closed-world}), this chapter develops the
comparison-completeness clause \cref{sp:closure}.
Building on the compactness boundary of \cref{ch:foundation}, it derives
the canonical product decomposition and diagonal-redundancy data used in
\cref{ch:rotation,ch:closure}.
\Cref{ch:foundation} isolated the Boolean compactness principle
governing the passage from finite-stage admissibility to global
admissibility.  In particular, \cref{thm:boundary,cor:admiss-inv,thm:stability}
show that admissible global objects exist precisely when no finite
family of excluded patterns already covers the universe.
The present chapter applies this compactness mechanism to the primitive
comparison data introduced in the Introduction.  Starting from a
comparison world
\[
	(U,\mathcal C),
\]
we analyze the intrinsic indistinguishability relations determined by
the comparison predicates and determine the exact condition under which
the comparison system admits a canonical global structure.
From the comparison predicates one obtains two intrinsic equivalence
relations recording indistinguishability of outgoing and incoming
comparison profiles.  These relations determine quotient sets
\[
	X_A := U/\alpha,
	\qquad
	X_B := U/\beta ,
\]
and hence a canonical factor map
\begin{equation}
	\label{eq:Theta}
	\Theta : U \to X_A \times X_B,
	\qquad
	\Theta(u) = ([u]_\alpha,[u]_\beta).
\end{equation}
The central structural question is therefore the following.
\medskip
\noindent
\emph{When does the intrinsic comparison structure determine the
	entire state space through this factor map?}
\medskip
The answer is a single closure principle: rectangular comparison
completeness, the chapter-level realization of
\cref{sp:closure}. We prove that this condition is equivalent to bijectivity
of \(\Theta\), and hence to the existence of a canonical product
decomposition
\[
	U \cong X_A \times X_B
\]
realizing the intrinsic congruences as coordinate equalities.
The proof proceeds in two complementary layers.
\begin{itemize}
	\item First, we analyze the comparison predicates directly and show
	      that rectangular completeness is equivalent to bijectivity of the
	      canonical factor map.
	\item Second, we internalize the same structure inside the Boolean
	      algebra generated by the comparison predicates.  The finite-coordinate
	      Boolean algebras form a refinement tower, and the compactness results
	      of \cref{ch:foundation} allow finite rectangle conditions to be lifted
	      to global splitting statements.
\end{itemize}
Under the standing principle, and in the chapter-local
distinguishability setup
\[
	\alpha \cap \beta = \Delta_U,
\]
these arguments yield equivalent formulations of rectangular
completeness in three different languages:
\begin{itemize}
	\item comparison structure;
	\item Boolean algebra factorization;
	\item ultrafilter splitting.
\end{itemize}
When these conditions hold, the intrinsic symmetry group
\[
	G := \Aut(U,\mathcal C)
\]
acts diagonally on the canonical product
\[
	X_A \times X_B,
\]
and the quotient data used later in the relational program are determined.
This establishes the canonical product decomposition promised in
\cref{thm:intro-structural-rigidity}.
% ============================================================
\section{Primitive comparison worlds}
% ============================================================
\begin{definition}[Comparison world (recalled)]
	As in \cref{def:intro-comparison-world}, a \emph{comparison world} is a pair \((U,\mathcal C)\), where
	\(U\neq\varnothing\) and \(\mathcal C\) is a family of maps
	\[
		c:U\times U\to\{0,1\}.
	\]
	In this manuscript such worlds are required to be intrinsic in the sense
	of \cref{def:intro-comparison-world}.
\end{definition}
\begin{remark}
	No symmetry, topology, product structure, or background geometry is
	assumed. Every structure appearing below is extracted from the
	comparison predicates themselves.
\end{remark}
\begin{definition}[Intrinsic symmetry group]
	\label{def:intrinsic-symmetry-group}
	Define
	\[
		G:=\Aut(U,\mathcal C)
		:=
		\bigl\{
		\phi\in\Bij(U):
		c(\phi(u),\phi(v))=c(u,v)
		\ \forall c\in\mathcal C,\ \forall u,v\in U
		\bigr\}.
	\]
\end{definition}
\begin{proposition}
	\(G\) is a group under composition.
\end{proposition}
\begin{proof}
	The identity map preserves every predicate, hence lies in \(G\).
	If \(\phi,\psi\in G\), then for every \(c\in\mathcal C\) and
	\(u,v\in U\),
	\[
		c((\phi\circ\psi)(u),(\phi\circ\psi)(v))
		=
		c(\phi(\psi(u)),\phi(\psi(v)))
		=
		c(\psi(u),\psi(v))
		=
		c(u,v),
	\]
	so \(\phi\circ\psi\in G\).
	If \(\phi\in G\), fix \(c\in\mathcal C\) and \(u,v\in U\). Write
	\(u=\phi(u')\) and \(v=\phi(v')\). Then
	\[
		c(\phi^{-1}(u),\phi^{-1}(v))
		=
		c(u',v')
		=
		c(\phi(u'),\phi(v'))
		=
		c(u,v).
	\]
	Hence \(\phi^{-1}\in G\).
\end{proof}
% ============================================================
\section{Intrinsic congruences}
% ============================================================
\begin{definition}[Left and right profiles]
	\label{def:left-right-profiles}
	For each \(u\in U\), define functions
	\[
		\begin{split}
			\mathsf L(u):\mathcal C\times U\to\{0,1\}, \\
			\mathsf L(u)(c,w):=c(u,w),
		\end{split}
	\]
	and
	\[
		\begin{split}
			\mathsf R(u):\mathcal C\times U\to\{0,1\}, \\
			\mathsf R(u)(c,w):=c(w,u).
		\end{split}
	\]
\end{definition}
\begin{definition}[Intrinsic congruences]
	\label{def:alpha-beta}
	Define relations \(\alpha,\beta\) on \(U\) by
	\[
		u\,\alpha\,v
		\iff
		\mathsf L(u)=\mathsf L(v),
		\qquad
		u\,\beta\,v
		\iff
		\mathsf R(u)=\mathsf R(v).
	\]
	Equivalently,
	\[
		u\,\alpha\,v
		\iff
		c(u,w)=c(v,w)\ \forall c\in\mathcal C,\ \forall w\in U,
	\]
	and
	\[
		u\,\beta\,v
		\iff
		c(w,u)=c(w,v)\ \forall c\in\mathcal C,\ \forall w\in U.
	\]
\end{definition}
\begin{lemma}
	\(\alpha\) and \(\beta\) are equivalence relations.
\end{lemma}
\begin{proof}
	Each is equality of functions in \(\{0,1\}^{\mathcal C\times U}\).
	Reflexivity, symmetry, and transitivity follow immediately.
\end{proof}
\begin{lemma}[Maximality among unary comparison invariants]
	Let \(\sim\) be any equivalence relation on \(U\) such that
	\[
		u\sim v
		\Longrightarrow
		c(u,w)=c(v,w)
		\qquad
		\forall c\in\mathcal C,\ \forall w\in U.
	\]
	Then \(\sim\subseteq\alpha\). Similarly, any equivalence relation
	preserving all right predicates is contained in \(\beta\).
\end{lemma}
\begin{proof}
	If \(u\sim v\), the stated hypothesis implies
	\(\mathsf L(u)=\mathsf L(v)\), hence \(u\,\alpha\,v\).
	Therefore every \(\sim\)-pair is an \(\alpha\)-pair.
	The right-handed statement is identical.
\end{proof}
\begin{lemma}[Invariance]
	\label{lem:invariance}
	\(\alpha\) and \(\beta\) are \(G\)-invariant.
\end{lemma}
\begin{proof}
	Suppose \(u\,\alpha\,v\), fix \(g\in G\), \(c\in\mathcal C\), and
	\(w'\in U\). Since \(g\) is bijective, there exists \(w\in U\) with
	\(w'=g(w)\). Then
	\[
		c(g(u),w')
		=
		c(g(u),g(w))
		=
		c(u,w)
		=
		c(v,w)
		=
		c(g(v),g(w))
		=
		c(g(v),w').
	\]
	Hence \(g(u)\,\alpha\,g(v)\). The proof for \(\beta\) is identical.
\end{proof}
% ============================================================
\section{The canonical factor map and its universal property}
\label{sec:Theta}
% ============================================================
Define
\[
	X_A:=U/\alpha,
	\qquad
	X_B:=U/\beta.
\]
\begin{definition}[Canonical factor map]
	\label{def:Theta}
	The canonical factor map is
	\[
		\Theta:U\to X_A\times X_B,
		\qquad
		\Theta(u):=([u]_\alpha,[u]_\beta).
	\]
\end{definition}
\begin{lemma}[Kernel characterization]
	\label{lem:kernel}
	For \(u,v\in U\),
	\[
		\Theta(u)=\Theta(v)
		\iff
		u\,\alpha\,v
		\text{ and }
		u\,\beta\,v.
	\]
	In particular,
	\[
		\Theta\text{ is injective}
		\iff
		\alpha\cap\beta=\Delta_U.
	\]
\end{lemma}
\begin{proof}
	By definition,
	\[
		\Theta(u)=\Theta(v)
		\iff
		[u]_\alpha=[v]_\alpha
		\text{ and }
		[u]_\beta=[v]_\beta,
	\]
	which is equivalent to \(u\,\alpha\,v\) and \(u\,\beta\,v\).
	The injectivity criterion is the special case of equality of points.
\end{proof}
\begin{definition}[Realizing product presentation]
	A \emph{realizing product presentation} is a bijection
	\[
		p:U\to Y_A\times Y_B
	\]
	such that
	\[
		u\,\alpha\,v
		\iff
		\pi_A(p(u))=\pi_A(p(v)),
	\]
	and
	\[
		u\,\beta\,v
		\iff
		\pi_B(p(u))=\pi_B(p(v)),
	\]
	where \(\pi_A,\pi_B\) denote the coordinate projections.
\end{definition}
\begin{proposition}[Terminality of the canonical factor map]
	\label{prop:terminality}
	Let
	\[
		p:U\to Y_A\times Y_B
	\]
	be a realizing product presentation. Then there exists a unique
	bijection
	\[
		\Phi:Y_A\times Y_B\to X_A\times X_B
	\]
	such that
	\[
		\Theta=\Phi\circ p.
	\]
\end{proposition}
\begin{proof}
	We define the coordinate maps separately.
	First define \(\phi_A:Y_A\to X_A\) by
	\[
		\phi_A(a):=[p^{-1}(a,b)]_\alpha,
	\]
	where \(b\in Y_B\) is arbitrary. We must show that this does not
	depend on the choice of \(b\). Let \(b_1,b_2\in Y_B\), and set
	\[
		u_i:=p^{-1}(a,b_i)\qquad (i=1,2).
	\]
	Then
	\[
		\pi_A(p(u_1))=\pi_A(p(u_2))=a.
	\]
	Since \(p\) realizes \(\alpha\), it follows that \(u_1\,\alpha\,u_2\),
	hence
	\[
		[p^{-1}(a,b_1)]_\alpha=[p^{-1}(a,b_2)]_\alpha.
	\]
	Thus \(\phi_A\) is well defined.
	Similarly define \(\phi_B:Y_B\to X_B\) by
	\[
		\phi_B(b):=[p^{-1}(a,b)]_\beta,
	\]
	where \(a\in Y_A\) is arbitrary. If \(a_1,a_2\in Y_A\), then with
	\(v_i:=p^{-1}(a_i,b)\) we have
	\[
		\pi_B(p(v_1))=\pi_B(p(v_2))=b.
	\]
	Since \(p\) realizes \(\beta\), we obtain \(v_1\,\beta\,v_2\), so
	\[
		[p^{-1}(a_1,b)]_\beta=[p^{-1}(a_2,b)]_\beta.
	\]
	Hence \(\phi_B\) is well defined.
	Now define
	\[
		\Phi:Y_A\times Y_B\to X_A\times X_B,
		\qquad
		\Phi(a,b):=(\phi_A(a),\phi_B(b)).
	\]
	This is well defined because \(\phi_A\) and \(\phi_B\) are.
	Let \(u\in U\), and write \(p(u)=(a,b)\). Then
	\[
		\Phi(p(u))
		=
		(\phi_A(a),\phi_B(b))
		=
		([u]_\alpha,[u]_\beta)
		=
		\Theta(u).
	\]
	Thus
	\[
		\Theta=\Phi\circ p.
	\]
	Uniqueness is immediate from bijectivity of \(p\): if \(\Phi'\) also
	satisfies \(\Theta=\Phi'\circ p\), then for every \((a,b)\in Y_A\times
	Y_B\),
	\[
		\Phi(a,b)=\Theta(p^{-1}(a,b))=\Phi'(a,b).
	\]
	It remains to prove that \(\Phi\) is bijective.
	\emph{Injectivity.}
	Assume
	\[
		\Phi(a_1,b_1)=\Phi(a_2,b_2).
	\]
	Set
	\[
		u_i:=p^{-1}(a_i,b_i)
		\qquad (i=1,2).
	\]
	Then
	\[
		\Theta(u_1)=\Phi(p(u_1))=\Phi(a_1,b_1)=\Phi(a_2,b_2)=\Phi(p(u_2))=\Theta(u_2).
	\]
	By \cref{lem:kernel},
	\[
		u_1\,\alpha\,u_2,
		\qquad
		u_1\,\beta\,u_2.
	\]
	Since \(p\) realizes \(\alpha\) and \(\beta\), this implies
	\[
		a_1=\pi_A(p(u_1))=\pi_A(p(u_2))=a_2,
		\qquad
		b_1=\pi_B(p(u_1))=\pi_B(p(u_2))=b_2.
	\]
	Hence \((a_1,b_1)=(a_2,b_2)\), so \(\Phi\) is injective.
	\emph{Surjectivity.}
	Let \((A,B)\in X_A\times X_B\). Choose \(u_A,u_B\in U\) such that
	\[
		[u_A]_\alpha=A,
		\qquad
		[u_B]_\beta=B.
	\]
	Write
	\[
		p(u_A)=(a_A,b_A),
		\qquad
		p(u_B)=(a_B,b_B).
	\]
	Set
	\[
		u:=p^{-1}(a_A,b_B).
	\]
	Then
	\[
		\pi_A(p(u))=a_A=\pi_A(p(u_A)).
	\]
	Since \(p\) realizes \(\alpha\), it follows that \(u\,\alpha\,u_A\),
	hence
	\[
		[u]_\alpha=[u_A]_\alpha=A.
	\]
	Likewise,
	\[
		\pi_B(p(u))=b_B=\pi_B(p(u_B)),
	\]
	so \(u\,\beta\,u_B\), hence
	\[
		[u]_\beta=[u_B]_\beta=B.
	\]
	Therefore
	\[
		\Theta(u)=([u]_\alpha,[u]_\beta)=(A,B).
	\]
	Since \(\Theta=\Phi\circ p\), we obtain
	\[
		(A,B)=\Theta(u)=\Phi(p(u)),
	\]
	and \(\Phi\) is surjective.
	Thus \(\Phi\) is bijective.
\end{proof}
% ============================================================
\section{Rectangular completeness}
\label{sec:rect}
% ============================================================
\begin{definition}[Rectangular completeness]
	\label{def:rect}
	The comparison world \((U,\mathcal C)\) is \emph{rectangularly complete}
	if for every \(A\in X_A\) and \(B\in X_B\) there exists a unique
	\(u\in U\) such that
	\[
		[u]_\alpha=A,
		\qquad
		[u]_\beta=B.
	\]
	Equivalently: for every \(u,v\in U\) there exists a unique \(x\in U\)
	such that \(x\,\alpha\,u\) and \(x\,\beta\,v\).
\end{definition}
\begin{remark}
	Rectangular completeness asserts that any left profile can be paired
	with any right profile, and that this pairing is unique.
\end{remark}
\subsection{A symmetry obstruction}
\label{sec:symmetry-obstruction}
\begin{definition}[Symmetric comparison world]
	A comparison world \((U,\mathcal C)\) is \emph{symmetric} if
	\[
		c(u,w)=c(w,u)
		\qquad
		\forall c\in\mathcal C,\ \forall u,w\in U.
	\]
\end{definition}
\begin{lemma}[Collapse of left and right profiles]
	\label{lem:LR-collapse}
	If \((U,\mathcal C)\) is symmetric, then
	\[
		\mathsf L(u)=\mathsf R(u)
		\qquad
		\forall u\in U.
	\]
	In particular,
	\[
		\alpha=\beta.
	\]
\end{lemma}
\begin{proof}
	For every \(u\in U\), \(c\in\mathcal C\), and \(w\in U\),
	\[
		\mathsf L(u)(c,w)=c(u,w)=c(w,u)=\mathsf R(u)(c,w).
	\]
	Hence \(\mathsf L(u)=\mathsf R(u)\). Therefore, for \(u,v\in U\),
	\[
		u\,\alpha\,v
		\iff
		\mathsf L(u)=\mathsf L(v)
		\iff
		\mathsf R(u)=\mathsf R(v)
		\iff
		u\,\beta\,v.
	\]
\end{proof}
\begin{proposition}[Symmetry obstructs rectangular completeness]
	\label{prop:symmetry-obstructs-rect}
	Assume \(|U|>1\) and \((U,\mathcal C)\) is symmetric. Then
	\((U,\mathcal C)\) is not rectangularly complete. Equivalently,
	\[
		\Theta:U\to X_A\times X_B
	\]
	is not bijective.
	More precisely, exactly one of the following occurs:
	\begin{enumerate}
		\item[\emph{(K)}] If \(\alpha\neq\Delta_U\), then \(\Theta\) is not injective.
		\item[\emph{(R)}] If \(\alpha=\Delta_U\), then \(\Theta\) is injective but not surjective.
	\end{enumerate}
\end{proposition}
\begin{proof}
	By \cref{lem:LR-collapse}, \(\alpha=\beta\).
	If \(\alpha\neq\Delta_U\), then
	\[
		\alpha\cap\beta=\alpha\neq\Delta_U,
	\]
	so \(\Theta\) is not injective by \cref{lem:kernel}.
	If \(\alpha=\Delta_U\), then also \(\beta=\Delta_U\), hence
	\[
		|X_A|=|U|,
		\qquad
		|X_B|=|U|.
	\]
	Therefore
	\[
		|X_A\times X_B|=|U|^2>|U|
	\]
	because \(|U|>1\). Since the image of \(\Theta\) has cardinality at most
	\(|U|\), \(\Theta\) cannot be surjective.
	Thus \(\Theta\) is never bijective.
\end{proof}
\begin{corollary}[Necessary asymmetry]
	Assume \(|U|>1\) and \((U,\mathcal C)\) is rectangularly complete.
	Then there exist \(c\in\mathcal C\) and \(u,w\in U\) such that
	\[
		c(u,w)\neq c(w,u).
	\]
\end{corollary}
\begin{proof}
	This is the contrapositive of
	\cref{prop:symmetry-obstructs-rect}.
\end{proof}
\subsection{Finite profile support and orbit separation}
\label{subsec:finite-profile-support}
The next definitions and lemmas record the finite-coordinate facts needed
for the intrinsicality consequence without routing through the later
Boolean factorization theorem.
\begin{definition}[Generated Boolean algebra]
	Let \(\B\) be the Boolean subalgebra of \(\mathcal P(U)\) generated by
	the sets
	\[
		L_{c,w}:=\{u\in U:\ c(u,w)=1\},
		\qquad
		R_{c,w}:=\{u\in U:\ c(w,u)=1\},
	\]
	for all \(c\in\mathcal C\) and \(w\in U\).
	Let \(\BA\subseteq\B\) be the Boolean subalgebra generated by the
	\(L_{c,w}\), and \(\BB\subseteq\B\) the Boolean subalgebra generated by
	the \(R_{c,w}\).
\end{definition}
\begin{definition}[Membership equivalence]
	For a family \(\mathcal F\subseteq\mathcal P(U)\), define a relation
	\(\equiv_{\mathcal F}\) on \(U\) by
	\[
		u\equiv_{\mathcal F} v
		\iff
		(\forall E\in\mathcal F)\,
		(u\in E \Leftrightarrow v\in E).
	\]
\end{definition}
\begin{lemma}[Classes as membership equivalence]
	\label{lem:classes-boolean}
	For \(u,v\in U\),
	\[
		u\,\alpha\,v \iff u\equiv_{\BA} v,
		\qquad
		u\,\beta\,v \iff u\equiv_{\BB} v.
	\]
\end{lemma}
\begin{proof}
	We prove the first equivalence; the second is identical.
	Assume \(u\,\alpha\,v\). Then by definition,
	\[
		c(u,w)=c(v,w)
		\qquad
		\forall c\in\mathcal C,\ \forall w\in U.
	\]
	Equivalently,
	\[
		u\in L_{c,w}
		\iff
		v\in L_{c,w}
		\qquad
		\forall c,w.
	\]
	Since \(\BA\) is the Boolean algebra generated by the \(L_{c,w}\),
	agreement on all generators implies agreement on every element of
	\(\BA\). Thus \(u\equiv_{\BA} v\).
	Conversely, if \(u\equiv_{\BA} v\), then in particular
	\[
		u\in L_{c,w}
		\iff
		v\in L_{c,w}
		\qquad
		\forall c,w,
	\]
	which means \(c(u,w)=c(v,w)\) for all \(c,w\). Hence \(u\,\alpha\,v\).
\end{proof}
\begin{definition}[Finite-coordinate subalgebras]
	Let \(\mathcal C_0\subseteq\mathcal C\) and \(W_0\subseteq U\) be finite.
	Define
	\[
		\BA(\mathcal C_0,W_0)
		:=
		\Bool\bigl(\{L_{c,w}:(c,w)\in \mathcal C_0\times W_0\}\bigr)
		\subseteq \BA,
	\]
	\[
		\BB(\mathcal C_0,W_0)
		:=
		\Bool\bigl(\{R_{c,w}:(c,w)\in \mathcal C_0\times W_0\}\bigr)
		\subseteq \BB,
	\]
	and
	\[
		\B(\mathcal C_0,W_0)
		:=
		\Bool\bigl(\BA(\mathcal C_0,W_0)\cup\BB(\mathcal C_0,W_0)\bigr)
		\subseteq \B.
	\]
\end{definition}
\begin{lemma}[Finite support]
	\label{lem:finite-support}
	Every element of \(\B\) belongs to \(\B(\mathcal C_0,W_0)\) for some
	finite \(\mathcal C_0\subseteq\mathcal C\) and finite \(W_0\subseteq U\).
	Likewise every element of \(\BA\) and of \(\BB\) belongs to a
	corresponding finite-coordinate subalgebra.
\end{lemma}
\begin{proof}
	By definition, \(\B\) is the Boolean algebra generated by the family
	\(\{L_{c,w},R_{c,w}\}_{(c,w)\in\mathcal C\times U}\). Every element of
	a generated Boolean algebra is obtained from finitely many generators
	using finitely many Boolean operations. Thus any given \(E\in\B\)
	depends on only finitely many pairs \((c,w)\). The same argument applies
	to \(\BA\) and \(\BB\).
\end{proof}
\begin{lemma}[Finite orbit separation under intrinsicality]
	\label{lem:orbit-separation-finite-stage}
	Let \((U,\mathcal C)\) be intrinsic, and let
	\(O_1,O_2\subseteq U\) be distinct \(G\)-orbits, where
	\(G=\Aut(U,\mathcal C)\). If representatives \(u\in O_1\) and
	\(v\in O_2\) are distinguished by the comparison data, in the sense that
	\(\mathsf L(u)\neq\mathsf L(v)\) or
	\(\mathsf R(u)\neq\mathsf R(v)\), then the distinction is witnessed by a
	single predicate coordinate. That is, there exist \(c\in\mathcal C\) and
	\(w\in U\) with
	\[
		c(u,w)\neq c(v,w)
	\]
	or with
	\[
		c(w,u)\neq c(w,v).
	\]
	Consequently the distinguishing cylinder lies in a finite-coordinate
	subalgebra \(\B(\mathcal C_0,W_0)\) with
	\(|\mathcal C_0|=|W_0|=1\), and no comparison-profile separation of
	orbits requires the completion \(\overline{\B}\).
\end{lemma}
\begin{proof}
	Suppose first that \(\mathsf L(u)\neq\mathsf L(v)\). By definition,
	\(\mathsf L(u)\) and \(\mathsf L(v)\) are functions
	\[
		\mathcal C\times U\to\{0,1\},
		\qquad
		\mathsf L(x)(c,w)=c(x,w).
	\]
	Inequality of functions means disagreement at some single point of the
	domain. Hence there is a pair \((c,w)\in\mathcal C\times U\) such that
	\[
		c(u,w)\neq c(v,w).
	\]
	Then one of \(u,v\) lies in \(L_{c,w}\) and the other does not, and
	\[
		L_{c,w}\in\B(\{c\},\{w\}).
	\]
	The right-profile case is identical: if \(\mathsf R(u)\neq\mathsf R(v)\),
	then for some \((c,w)\) one has \(c(w,u)\neq c(w,v)\), and the cylinder
	\(R_{c,w}\in\B(\{c\},\{w\})\) witnesses the distinction.

	Intrinsicality makes this single-coordinate witness internal rather than
	an external label. If \(g\in G\), then
	\[
		c(gu,gw)=c(u,w)
		\qquad\text{and}\qquad
		c(gv,gw)=c(v,w),
	\]
	so any left-coordinate distinction is transported equivariantly to the
	diagonal translate \((gu,gv)\) by the single transported coordinate
	\((c,gw)\); the right-coordinate case is the same. Thus a comparison
	distinction cannot first appear as an infinite or completed Boolean
	combination. If two profiles differ, they differ at one coordinate; if
	they agree at every coordinate, they are the same profile.
\end{proof}
\begin{remark}[Why no limit-only separation exists]
	\label{rem:no-limit-only-separation}
	\Cref{lem:orbit-separation-finite-stage} is the point at which
	intrinsicality excludes finitely unwitnessed profile classes. Inequality
	of profile functions is disagreement at a single coordinate; there is no
	notion of two profiles that agree at every finite coordinate but differ
	in the limit, because functions equal at every coordinate are equal. A
	distinction surviving only in a completion
	\(\overline{\B}\setminus\B\) would require the individual predicate
	coordinates not to register the distinction while an infinite Boolean
	combination does. But each predicate is a \(\{0,1\}\)-valued coordinate,
	and by intrinsicality each coordinate is transported by the symmetries it
	helps constitute. The completion adjoins no new comparison separations;
	it adjoins only infinite Boolean combinations of separations already
	witnessed coordinatewise.
\end{remark}
% ============================================================
\section{Main classification theorem}
\label{sec:classification}
% ============================================================
\begin{theorem}[Exact classification]
	\label{thm:classification}
	The following are equivalent:
	\begin{enumerate}
		\item[\emph{(i)}] \((U,\mathcal C)\) is rectangularly complete.
		\item[\emph{(ii)}] \(\Theta:U\to X_A\times X_B\) is bijective.
		\item[\emph{(iii)}] There exists a realizing product presentation.
	\end{enumerate}
	Moreover, when these conditions hold, the identification
	\[
		U\cong X_A\times X_B
	\]
	is canonical, namely \(\Theta\), and any realizing presentation factors
	uniquely through \(\Theta\) by \cref{prop:terminality}.
\end{theorem}
\begin{proof}
	\emph{(i)\(\Rightarrow\)(ii).}
	Rectangular completeness is exactly the statement that every pair
	\((A,B)\in X_A\times X_B\) has a unique preimage under \(\Theta\).
	\emph{(ii)\(\Rightarrow\)(i).}
	If \(\Theta\) is bijective, then for every \((A,B)\in X_A\times X_B\)
	there is a unique
	\[
		u=\Theta^{-1}(A,B)
	\]
	with \([u]_\alpha=A\) and \([u]_\beta=B\).
	\emph{(ii)\(\Rightarrow\)(iii).}
	Take \(Y_A=X_A\), \(Y_B=X_B\), and \(p=\Theta\).
	\emph{(iii)\(\Rightarrow\)(ii).}
	Apply \cref{prop:terminality}. If \(p\) is a realizing product
	presentation, then there exists a bijection
	\[
		\Phi:Y_A\times Y_B\to X_A\times X_B
	\]
	with
	\[
		\Theta=\Phi\circ p.
	\]
	Since \(p\) is bijective, \(\Theta\) is bijective.
\end{proof}
\begin{lemma}[Conservative-completion dichotomy]
	\label{lem:conservative-completion-dichotomy}
	Let \((U,\mathcal C)\) be intrinsic and locally distinguishable, and let
	\((A,B)\in X_A\times X_B\) be a profile pair with no \(u\in U\)
	satisfying \([u]_\alpha=A\) and \([u]_\beta=B\). Choose representatives
	\(a\in A\) and \(b\in B\), and write the class-determined values
	\[
		A(c,w):=c(a,w),
		\qquad
		B(c,w):=c(w,b),
	\]
	which are independent of the chosen representatives by
	\cref{def:alpha-beta}. Form the one-point extension
	\[
		U^+:=U\cup\{u_{AB}\},
	\]
	where \(u_{AB}\) is a formal element assigned the comparisons of profile
	\((A,B)\): for each \(c\in\mathcal C\) and \(w\in U\), set
	\[
		c(u_{AB},w):=A(c,w),
		\qquad
		c(w,u_{AB}):=B(c,w),
	\]
	and choose \(c(u_{AB},u_{AB})\) whenever the same admissibility conditions
	allow a diagonal value. Then exactly one of the following holds.
	\begin{enumerate}[label=\textup{(\alph*)}]
		\item\label{case:internal-exclusion}
		      \emph{Internal exclusion.} The assignment is inconsistent with the
		      comparison data: some \(c\in\mathcal C\), together with the
		      structural conditions defining admissible profiles in
		      \cref{def:left-right-profiles,def:alpha-beta}, forbids a state of
		      profile \((A,B)\). Then \((A,B)\) is not an admissible profile pair,
		      and no obstruction to surjectivity arises from it.
		\item\label{case:conservative-extension}
		      \emph{Conservative extension.} The assignment is consistent, and
		      the inclusion \(U\hookrightarrow U^+\) preserves the finitely
		      supported comparison algebra and automorphism-invariant profile data:
		      restriction from \(U^+\) to \(U\) is an isomorphism on every
		      finite-coordinate subalgebra supported in \(U\), and each
		      \(g\in\Stab_G(A,B)\) extends to an automorphism of \(U^+\) fixing
		      \(u_{AB}\). Equivalently, the full \(G\)-action extends after
		      adjoining the \(G\)-orbit of the formal profile pair. In this case
		      the presence or absence of \(u_{AB}\) is not registered by any
		      intrinsic predicate, invariant, or finite comparison context
		      supported in \(U\); the two worlds are intrinsically
		      indistinguishable with respect to the pair \((A,B)\).
	\end{enumerate}
	There is no third case: an extension that is neither internally excluded
	nor conservative is impossible.
\end{lemma}
\begin{proof}
	The alternatives are mutually exclusive by definition: in
	\ref{case:internal-exclusion} no consistent assignment exists, while
	\ref{case:conservative-extension} begins with a consistent assignment.
	It remains to show exhaustiveness. Suppose
	\ref{case:internal-exclusion} fails, so the formal profile assignment is
	consistent. We prove that \ref{case:conservative-extension} holds.

	\emph{Finite-coordinate algebra is preserved.} By
	\cref{lem:finite-support}, every element of \(\B\) has finite support,
	depending on finitely many pairs \((c,w)\) with \(w\in U\). Adjoining
	\(u_{AB}\) introduces evaluations at the new point, but the values
	\(c(u_{AB},w)\) and \(c(w,u_{AB})\) are fixed by the existing profile
	classes \(A\) and \(B\). Hence no finite-coordinate subalgebra supported
	in \(U\) acquires a new value not already dictated by \(U\)'s comparison
	data, and restriction from \(U^+\) to \(U\) is a Boolean isomorphism on
	each such finite-coordinate subalgebra.

	\emph{No new profile distinction is created on old states.} Suppose, for
	contradiction, that adjoining \(u_{AB}\) refined a profile class on
	\(U\). By \cref{lem:orbit-separation-finite-stage}, any such separation
	is witnessed at a single predicate coordinate. If the coordinate is
	\((c,w)\) with \(w\in U\), the disagreement already existed in \(U\), a
	contradiction. Thus the only possible new coordinate uses
	\(w=u_{AB}\). If \(u\,\alpha\,v\) in \(U\), then
	\[
		c(u,u_{AB})=B(c,u)=c(u,b)=c(v,b)=B(c,v)=c(v,u_{AB}),
	\]
	so the new coordinate does not split an old \(\alpha\)-class. Similarly,
	if \(u\,\beta\,v\) in \(U\), then
	\[
		c(u_{AB},u)=A(c,u)=c(a,u)=c(a,v)=A(c,v)=c(u_{AB},v),
	\]
	so it does not split an old \(\beta\)-class. Therefore the profile
	quotients of \(U^+\) restrict to the original profile quotients on
	\(U\).

	\emph{Automorphisms extend on the profile stabilizer.} Let
	\(g\in\Stab_G(A,B)\). Define \(g^+\) on \(U^+\) by
	\(g^+|_U=g\) and \(g^+(u_{AB})=u_{AB}\). Since \(g\) preserves all
	predicates on \(U\), it remains only to check pairs involving
	\(u_{AB}\). For \(w\in U\), the condition \(gA=A\) gives
	\[
		c(g^+u_{AB},gw)=c(u_{AB},gw)=A(c,gw)=A(c,w)=c(u_{AB},w),
	\]
	and \(gB=B\) gives
	\[
		c(gw,g^+u_{AB})=c(gw,u_{AB})=B(c,gw)=B(c,w)=c(w,u_{AB}).
	\]
	The diagonal value is fixed by consistency of the assigned profile.
	Thus \(g^+\in\Aut(U^+,\mathcal C)\). For a general \(g\in G\), the same
	calculation sends \(u_{AB}\) to the formal point of profile
	\((gA,gB)\), so the full action extends after adjoining the orbit of the
	missing profile pair.

	These three facts are precisely the conservative-extension alternative,
	so if internal exclusion fails, conservative extension holds. Hence the
	dichotomy is exhaustive.
\end{proof}
\begin{corollary}[Closure forbids conservative omission]
	\label{cor:closure-forbids-conservative-omission}
	In a world closed under intrinsic description, the conservative-extension
	case of \cref{lem:conservative-completion-dichotomy} cannot obtain for
	any unrealized admissible profile pair. For in that case the omission of
	\(u_{AB}\) is maintained by no intrinsic predicate, invariant, or finite
	comparison context supported in \(U\); by the lemma it is intrinsically
	undetectable. Such an omission therefore depends on a selection external
	to the comparison structure. Closure, as expressed by intrinsicality in
	\cref{def:intro-comparison-world} together with the prohibition on
	external scaffolding in \cref{sp:intrinsic}, admits no externally
	maintained omission. Hence every admissible profile pair is either
	internally excluded or already realized; equivalently, no unrealized
	admissible rectangular cell remains.
\end{corollary}
\begin{theorem}[Rectangular completeness under intrinsicality]
	\label{thm:rectangular-completeness-intrinsicality}
	Let \((U,\mathcal C)\) be intrinsic (\cref{def:intro-comparison-world})
	with local distinguishability
	\[
		\alpha\cap\beta=\Delta_U.
	\]
	Then \(\Theta\) is bijective; the world is rectangularly complete.
\end{theorem}
\begin{proof}
	Injectivity is \cref{lem:kernel} under
	\(\alpha\cap\beta=\Delta_U\). For surjectivity, suppose
	\((A,B)\in X_A\times X_B\) is an admissible profile pair unrealized in
	\(U\). By \cref{lem:conservative-completion-dichotomy}, the one-point
	completion at \((A,B)\) is either internally excluded or conservative.
	Internal exclusion contradicts admissibility. If the completion is
	conservative, \cref{cor:closure-forbids-conservative-omission} applies:
	a closed world admits no conservative omission, so \((A,B)\) must be
	realized. Either way no unrealized admissible cell survives. Therefore
	\(\Theta\) is surjective, hence bijective.
\end{proof}
\begin{remark}
	Equivalently, one may realize the missing cells by finite induction. By
	\cref{lem:conservative-completion-dichotomy}, each conservative adjunction
	leaves the finitely supported profile algebra of \cref{lem:finite-support}
	unchanged and reduces the count of unrealized admissible pairs; by
	\cref{lem:orbit-separation-finite-stage}, no separation requires a
	completion. The induction and the dichotomy are the same argument viewed
	from opposite sides: the engine of surjectivity is closure's prohibition
	on intrinsically undetectable omission, not the bookkeeping of the
	iteration.
\end{remark}
\begin{remark}
	\label{rem:intrinsic-rect-not-universal}
	This does not make rectangular completeness universal. Non-intrinsic
	worlds, and intrinsic worlds failing \(\alpha\cap\beta=\Delta_U\), may
	fail it; the symmetric worlds of \cref{prop:symmetry-obstructs-rect}
	remain genuine open systems. The theorem states that an intrinsic,
	locally distinguishing world is closed: closedness is a consequence of
	the primitive being a genuine comparison, not an added assumption.
\end{remark}
% ============================================================
\section{Failure modes}
\label{sec:failure}
% ============================================================
\begin{definition}[Kernel obstruction]
	Define
	\[
		K_{\alpha\beta}:=\alpha\cap\beta.
	\]
\end{definition}
\begin{proposition}[Injectivity obstruction]
	\label{prop:inj-obstruction}
	\(\Theta\) is injective if and only if \(K_{\alpha\beta}=\Delta_U\).
\end{proposition}
\begin{proof}
	This is exactly \cref{lem:kernel}.
\end{proof}
\begin{definition}[Rectangular deficiency]
	We say that \((U,\mathcal C)\) has \emph{rectangular deficiency} if
	\(K_{\alpha\beta}=\Delta_U\) but \(\Theta\) is not surjective.
	Equivalently, there exist \(A\in X_A\) and \(B\in X_B\) such that no
	\(u\in U\) satisfies
	\[
		([u]_\alpha,[u]_\beta)=(A,B).
	\]
\end{definition}
\begin{theorem}[Exhaustive trichotomy]
	Exactly one of the following holds:
	\begin{enumerate}
		\item[\emph{(D)}] \(\Theta\) is bijective.
		\item[\emph{(K)}] \(K_{\alpha\beta}\neq\Delta_U\).
		\item[\emph{(R)}] \(K_{\alpha\beta}=\Delta_U\) but \(\Theta\) is not surjective.
	\end{enumerate}
\end{theorem}
\begin{proof}
	Either \(\Theta\) is bijective or it is not. If it is not bijective,
	then either injectivity fails or surjectivity fails. By
	\cref{prop:inj-obstruction}, injectivity fails exactly when
	\(K_{\alpha\beta}\neq\Delta_U\). If injectivity holds but surjectivity fails, then
	by definition we are in the rectangle-deficient case.
	The three cases are mutually exclusive and exhaustive.
\end{proof}
\begin{proposition}[Minimal counterexamples]
	There exist comparison worlds realizing each obstruction:
	\begin{enumerate}
		\item a kernel-obstructed example with \(|U|=2\);
		\item a rectangle-deficient example with \(|U|=3\) and \(K_{\alpha\beta}=\Delta_U\).
	\end{enumerate}
\end{proposition}
\begin{proof}
	For kernel obstruction, let \(U=\{a,b\}\) and let
	\(\mathcal C=\{c\}\) with \(c(u,v)=0\) for all pairs.
	Then all left and right profiles are identical, so
	\[
		\alpha=\beta=U\times U,
		\qquad
		K_{\alpha\beta}\neq\Delta_U.
	\]
	For rectangle deficiency, let \(U=\{x,y,z\}\) and let
	\(\mathcal C=\{c\}\), where \(c(u,v)=1\) if and only if \(u=v\).
	Then both left and right profiles separate points, so
	\[
		\alpha=\beta=\Delta_U,
		\qquad
		K_{\alpha\beta}=\Delta_U.
	\]
	Hence \(|X_A|=|X_B|=3\), while \(|U|=3\). Therefore
	\[
		|X_A\times X_B|=9>|U|,
	\]
	so \(\Theta\) cannot be surjective.
\end{proof}
% ============================================================
\section{Canonical diagonal symmetry once decomposable}
\label{sec:diagonal}
% ============================================================
Assume from now on that the equivalent conditions of
\cref{thm:classification} hold, so that \(\Theta\) is bijective.
\begin{definition}[Induced factor actions]
	\label{def:induced-actions}
	For \(g\in G\) define actions on \(X_A\) and \(X_B\) by
	\[
		g\cdot [u]_\alpha := [g(u)]_\alpha,
		\qquad
		g\cdot [u]_\beta := [g(u)]_\beta.
	\]
\end{definition}
\begin{lemma}[Well-definedness]
	\label{lem:well-defined-actions}
	The actions in \cref{def:induced-actions} are well defined.
\end{lemma}
\begin{proof}
	If \([u]_\alpha=[v]_\alpha\), then \(u\,\alpha\,v\). By
	\cref{lem:invariance}, \(g(u)\,\alpha\,g(v)\), hence
	\([g(u)]_\alpha=[g(v)]_\alpha\). The proof for \(X_B\) is identical.
\end{proof}
\begin{theorem}[Diagonal action theorem]
	\label{thm:diagonal}
	Under the identification \(U\cong X_A\times X_B\) via \(\Theta\),
	the group \(G=\Aut(U,\mathcal C)\) acts diagonally:
	\[
		\Theta(g(u))=g\cdot\Theta(u)
		\qquad
		\forall g\in G,\ \forall u\in U,
	\]
	that is,
	\[
		g\cdot(A,B):=(g\cdot A,\ g\cdot B).
	\]
\end{theorem}
\begin{proof}
	By definition,
	\[
		\Theta(g(u))
		=
		([g(u)]_\alpha,[g(u)]_\beta)
		=
		(g\cdot[u]_\alpha,\ g\cdot[u]_\beta)
		=
		g\cdot\Theta(u).
	\]
\end{proof}
\begin{corollary}[Canonical diagonal redundancy data]
	\label{cor:terminal-data}
	Under rectangular completeness, the canonical data
	\[
		X_A:=U/\alpha,
		\qquad
		X_B:=U/\beta,
		\qquad
		G:=\Aut(U,\mathcal C)
	\]
	determine a canonical diagonal action on
	\[
		X:=X_A\times X_B
	\]
	and hence a canonical quotient
	\[
		\Phys:=X/G.
	\]
\end{corollary}
\begin{proof}
	This is immediate from \cref{thm:classification,thm:diagonal}.
\end{proof}
% ============================================================
\section{Boolean algebra internalization}
\label{sec:boolean}
% ============================================================
We now reformulate the structure of
\cref{thm:classification,thm:diagonal,cor:terminal-data} inside the
Boolean algebra
generated by the primitive comparison predicates.
% ------------------------------------------------------------
\subsection{The comparison Boolean algebra}
% ------------------------------------------------------------
The comparison Boolean algebra, membership equivalence, finite-coordinate
subalgebras, and finite-support lemmas were introduced in
\cref{subsec:finite-profile-support}. We use those definitions throughout
the Boolean internalization below.
\begin{remark}
	All atomic arguments below are carried out inside finite-coordinate
	subalgebras. No global atom of the full algebra \(\B\) is invoked.
\end{remark}
% ------------------------------------------------------------
\subsection{Finite-predicate refinement tower and compactness}
\label{subsec:refinement-compactness}
% ------------------------------------------------------------
We now make explicit the refinement tower to which
\cref{ch:foundation} applies. This is the point at which the finite
Boolean structure of the comparison predicates is connected to the
global admissibility mechanism isolated in \cref{ch:foundation}.
Fix an increasing exhaustion
\[
	F_0\subseteq F_1\subseteq\cdots\subseteq \mathcal C\times U,
	\qquad
	\bigcup_{n\in\mathbb N}F_n=\mathcal C\times U.
\]
Define
\[
	\B^{(n)}
	:=
	\Bool\bigl(\{L_{c,w},R_{c,w}:(c,w)\in F_n\}\bigr)
	\subseteq \mathcal P(U),
\]
\[
	\BA^{(n)}
	:=
	\Bool\bigl(\{L_{c,w}:(c,w)\in F_n\}\bigr),
	\qquad
	\BB^{(n)}
	:=
	\Bool\bigl(\{R_{c,w}:(c,w)\in F_n\}\bigr).
\]
Then
\[
	\B^{(n)}\subseteq \B^{(n+1)},
	\qquad
	\BA^{(n)}\subseteq \BA^{(n+1)},
	\qquad
	\BB^{(n)}\subseteq \BB^{(n+1)},
\]
and
\[
	\B=\bigcup_{n\in\mathbb N}\B^{(n)},
	\qquad
	\BA=\bigcup_{n\in\mathbb N}\BA^{(n)},
	\qquad
	\BB=\bigcup_{n\in\mathbb N}\BB^{(n)}.
\]
For each \(n\), let \(\mathrm{At}(\BA^{(n)})\) and
\(\mathrm{At}(\BB^{(n)})\) denote the atoms of these finite Boolean
algebras. Define \(\J^{(n)}\subseteq \B^{(n)}\) to be the ideal
generated by the empty rectangles
\[
	A\cap B,
	\qquad
	A\in\mathrm{At}(\BA^{(n)}),
	\quad
	B\in\mathrm{At}(\BB^{(n)}),
	\quad
	A\cap B=\varnothing.
\]
\begin{lemma}[Compatibility of the finite rectangle ideals]
	For every \(n\),
	\[
		\J^{(n)}=\J^{(n+1)}\cap \B^{(n)}.
	\]
\end{lemma}
\begin{proof}
	Every atom of \(\BA^{(n+1)}\) is contained in a unique atom of
	\(\BA^{(n)}\), and similarly for \(\BB^{(n+1)}\) and \(\BB^{(n)}\).
	Thus an empty rectangle at stage \(n+1\) restricts to an empty
	rectangle at stage \(n\). Conversely, an empty rectangle at stage \(n\)
	is the finite union of all stage-\((n+1)\) rectangles lying over it,
	and each of those is empty. The generated ideals therefore agree under
	restriction.
\end{proof}
Set
\[
	\J_\infty:=\left\langle \bigcup_{n\in\mathbb N}\J^{(n)}\right\rangle_{\B}.
\]
By \cref{thm:boundary,cor:admiss-inv,thm:stability}, the tower
\[
	\bigl(U,(\B^{(n)})_{n\in\mathbb N},(\J^{(n)})_{n\in\mathbb N}\bigr)
\]
admits a global admissible ultrafilter exactly when no finite family of
excluded finite-coordinate rectangles covers \(U\), and admissible
global ultrafilters are exactly coherent admissible towers of stage
ultrafilters:
\[
	\UF(\B;\J_\infty)
	\cong
	\varprojlim_n \UF(\B^{(n)};\J^{(n)}).
\]
This is the compactness input used in the reverse implication of
\cref{thm:boolean-complete}.
\begin{theorem}[Compactness from intrinsicality]
	\label{thm:compactness-from-intrinsicality}
	For an intrinsic comparison world (\cref{def:intro-comparison-world}),
	the finite-coordinate refinement tower has limit equal to the direct
	union
	\[
		\B=\bigcup_{n\in\mathbb N}\B^{(n)},
	\]
	and every comparison profile class separated by intrinsic comparison data
	is decided at a finite stage. Consequently the finite-to-global boundary
	holds without separate hypothesis: global admissibility fails if and
	only if finitely many stage-exclusions cover \(U\), as in
	\cref{thm:boundary}.
\end{theorem}
\begin{proof}
	The equality \(\B=\bigcup_n\B^{(n)}\) is exactly the finite-support
	content of \cref{lem:finite-support} applied to the chosen exhaustion of
	\(\mathcal C\times U\). Intrinsicality makes the profile information
	internal to the symmetry structure: each predicate is preserved by every
	element of \(G=\Aut(U,\mathcal C)\), and the family of left and right
	cylinders is transported within the comparison-generated Boolean algebra.

	Suppose a comparison profile class required a completion element outside
	the finite union tower. Then it would determine a profile distinction not
	witnessed by any finite-coordinate subalgebra. This is impossible by
	\cref{lem:orbit-separation-finite-stage}: any distinction between left or
	right profile functions is witnessed by a single predicate coordinate.
	That coordinate lies in a finite-coordinate subalgebra, and
	\cref{lem:finite-support} then places every Boolean combination of such
	witnesses in some \(\B^{(n)}\). Thus no profile class is finitely
	unwitnessed, the limit requires no additional completion stage, and the
	admissibility boundary is precisely the quantifier-reversal boundary of
	\cref{thm:boundary}.
\end{proof}
% ------------------------------------------------------------
\subsection{Rectangular completeness as a Boolean factorization property}
% ------------------------------------------------------------
From this point onward in the Boolean part of the chapter, under the
comparison-completeness clause \cref{sp:closure}, we work in the local
distinguishability setup
\[
	\alpha\cap\beta=\Delta_U.
\]
\begin{definition}[Rectangular factorization property]
	We say that \(\B\) has the \emph{rectangular factorization property} if:
	\begin{enumerate}
		\item[\emph{(F1)}]
		      for every nonempty \(A\in\BA\) and nonempty \(B\in\BB\),
		      \[
			      A\cap B\neq\varnothing;
		      \]
		\item[\emph{(F2)}]
		      every \(E\in\B\) is a finite union of sets of the form \(A\cap B\),
		      with \(A\in\BA\) and \(B\in\BB\).
	\end{enumerate}
\end{definition}
\begin{theorem}[Boolean internalization of rectangular completeness]
	\label{thm:boolean-complete}
	In the local distinguishability setup
	\[
		\alpha\cap\beta=\Delta_U.
	\]
	Then rectangular completeness holds if and only if \(\B\) satisfies the
	rectangular factorization property.
\end{theorem}
\begin{proof}
	\emph{(\(\Rightarrow\)).}
	Assume rectangular completeness. By \cref{thm:classification},
	\(\Theta:U\to X_A\times X_B\) is bijective. Identify \(U\) with
	\(X_A\times X_B\) through \(\Theta\).
	If \(A\in\BA\), then by \cref{lem:classes-boolean} membership in \(A\)
	depends only on the \(\alpha\)-class, hence only on the first
	coordinate. Therefore there exists a subset \(A_A\subseteq X_A\) such
	that
	\[
		A=A_A\times X_B.
	\]
	Similarly, every \(B\in\BB\) has the form
	\[
		B=X_A\times B_B
	\]
	for some \(B_B\subseteq X_B\).
	Now if \(A\in\BA\) and \(B\in\BB\) are nonempty, then
	\[
		A=A_A\times X_B,
		\qquad
		B=X_A\times B_B
	\]
	with \(A_A\neq\varnothing\) and \(B_B\neq\varnothing\), so
	\[
		A\cap B=(A_A\times X_B)\cap(X_A\times B_B)=A_A\times B_B\neq\varnothing.
	\]
	Thus \emph{(F1)} holds.
	Moreover, \(\B\) is generated by the left and right coordinate
	cylinders, and every finite Boolean combination of such cylinders is a
	finite union of rectangles \(A_A\times B_B\). Transporting back through
	\(\Theta\) yields \emph{(F2)}.
	\medskip
	\emph{(\(\Leftarrow\)).}
	Assume \emph{(F1)} and \emph{(F2)}.
	We first prove surjectivity of \(\Theta\). Fix \(A\in X_A\) and
	\(B\in X_B\), and choose representatives \(u_A,u_B\in U\) such that
	\[
		[u_A]_\alpha=A,
		\qquad
		[u_B]_\beta=B.
	\]
	For each \(n\), let \(A_n\in\mathrm{At}(\BA^{(n)})\) be the unique atom
	containing \(u_A\), and let \(B_n\in\mathrm{At}(\BB^{(n)})\) be the
	unique atom containing \(u_B\). Since both are nonempty, \emph{(F1)}
	implies
	\[
		A_n\cap B_n\neq\varnothing
		\qquad
		\text{for every }n.
	\]
	Choose \(x_n\in A_n\cap B_n\), and let \(\mathfrak u_n\) be the
	principal ultrafilter of \(\B^{(n)}\) determined by \(x_n\). Because
	\(x_n\in A_n\cap B_n\) and the ideal \(\J^{(n)}\) is generated by empty
	rectangles, one has
	\[
		\mathfrak u_n\in \UF(\B^{(n)};\J^{(n)}).
	\]
	Moreover, the atoms \(A_{n+1}\subseteq A_n\) and \(B_{n+1}\subseteq B_n\)
	refine coherently, so the stage ultrafilters may be chosen coherently.
	By \cref{cor:admiss-inv}, there exists
	\[
		\mathfrak u\in \UF(\B;\J_\infty)
	\]
	whose restriction to \(\B^{(n)}\) is \(\mathfrak u_n\) for every \(n\).
	By construction, \(\mathfrak u\) contains every set \(E\in\BA\)
	containing \(u_A\) and every set \(F\in\BB\) containing \(u_B\). Hence
	the complete left-membership data encoded by \(\mathfrak u\) agree with
	those of \(u_A\), and the complete right-membership data agree with
	those of \(u_B\). Equivalently, \(\mathfrak u\) determines a state
	\(u\in U\) satisfying
	\[
		[u]_\alpha=[u_A]_\alpha=A,
		\qquad
		[u]_\beta=[u_B]_\beta=B.
	\]
	Thus \(\Theta\) is surjective.
	We now prove injectivity. Since
	\[
		\alpha\cap\beta=\Delta_U,
	\]
	injectivity of \(\Theta\) is immediate from \cref{lem:kernel}.
	Therefore \(\Theta\) is bijective. By \cref{thm:classification},
	rectangular completeness follows.
\end{proof}
% ------------------------------------------------------------
\subsection{Rectangle algebras and ultrafilter splitting}
% ------------------------------------------------------------
We now reformulate the same condition in algebraic and ultrafilter
language.
\begin{definition}[Ultrafilter]
	Let \(\mathcal A\) be a Boolean algebra. An \emph{ultrafilter} on
	\(\mathcal A\) is a subset \(\mathfrak u\subseteq\mathcal A\) such that:
	\begin{enumerate}
		\item \(1\in\mathfrak u\) and \(0\notin\mathfrak u\);
		\item \(a,b\in\mathfrak u\Rightarrow a\wedge b\in\mathfrak u\);
		\item \(a\in\mathfrak u\) and \(a\le b\Rightarrow b\in\mathfrak u\);
		\item for each \(a\in\mathcal A\), exactly one of \(a\) and \(\neg a\)
		      lies in \(\mathfrak u\).
	\end{enumerate}
\end{definition}
\begin{definition}[Marginals]
	For \(\mathfrak u\in\UF(\B)\), define
	\[
		\mathfrak u_A:=\mathfrak u\cap\BA\in\UF(\BA),
		\qquad
		\mathfrak u_B:=\mathfrak u\cap\BB\in\UF(\BB).
	\]
\end{definition}
\begin{lemma}
	\label{lem:marginals}
	If \(\mathfrak u\in\UF(\B)\), then \(\mathfrak u_A\in\UF(\BA)\) and
	\(\mathfrak u_B\in\UF(\BB)\).
\end{lemma}
\begin{proof}
	Restriction of an ultrafilter to a Boolean subalgebra preserves all
	ultrafilter axioms.
\end{proof}
\begin{definition}[Rectangle algebra]
	Let \(\mathcal A,\mathcal D\) be Boolean algebras. A \emph{rectangle
		algebra for \((\mathcal A,\mathcal D)\)} is a Boolean algebra
	\(\mathcal T\) equipped with homomorphisms
	\[
		\iota_A:\mathcal A\to\mathcal T,
		\qquad
		\iota_D:\mathcal D\to\mathcal T
	\]
	such that:
	\begin{enumerate}
		\item the images \(\iota_A(\mathcal A)\cup\iota_D(\mathcal D)\)
		      generate \(\mathcal T\);
		\item for any Boolean algebra \(\mathcal E\) and homomorphisms
		      \(f:\mathcal A\to\mathcal E\), \(g:\mathcal D\to\mathcal E\), there
		      exists a unique homomorphism \(h:\mathcal T\to\mathcal E\) such that
		      \[
			      h\circ\iota_A=f,
			      \qquad
			      h\circ\iota_D=g.
		      \]
	\end{enumerate}
	When such an algebra exists, we write
	\[
		\mathcal T=\mathcal A\otimes\mathcal D.
	\]
\end{definition}
\begin{proposition}[Existence inside \(\mathcal P(U)\)]
	\label{prop:tensor-exists}
	Let \(\mathcal T\) be the Boolean subalgebra of \(\mathcal P(U)\)
	generated by all intersections \(A\cap B\) with \(A\in\BA\) and
	\(B\in\BB\). Then \(\mathcal T\) is a rectangle algebra for
	\((\BA,\BB)\).
\end{proposition}
\begin{proof}
	Set
	\[
		\mathcal T:=\Bool\bigl(\{A\cap B:\ A\in\BA,\ B\in\BB\}\bigr)
		\subseteq\mathcal P(U).
	\]
	Define
	\[
		\iota_A(A):=A,
		\qquad
		\iota_D(B):=B.
	\]
	The images generate \(\mathcal T\) by construction.
	Now let \(f:\BA\to\mathcal E\) and \(g:\BB\to\mathcal E\) be Boolean
	homomorphisms into a Boolean algebra \(\mathcal E\). On rectangle
	generators define
	\[
		h(A\cap B):=f(A)\wedge g(B).
	\]
	Because \(\mathcal T\) is generated by these intersections and Boolean
	operations, \(h\) extends uniquely to a homomorphism
	\[
		h:\mathcal T\to\mathcal E.
	\]
	This gives the universal property.
\end{proof}
\begin{theorem}[Ultrafilter splitting characterization]
	\label{thm:ultrafilter-splitting}
	In the local distinguishability setup
	\[
		\alpha\cap\beta=\Delta_U.
	\]
	The following are equivalent:
	\begin{enumerate}
		\item[\emph{(a)}] Rectangular completeness holds.
		\item[\emph{(b)}] \(\B\) satisfies the rectangular factorization property.
		\item[\emph{(c)}] \(\B\cong\BA\otimes\BB\).
		\item[\emph{(d)}] The marginal map
		      \[
			      \Sigma:\UF(\B)\to\UF(\BA)\times\UF(\BB),
			      \qquad
			      \Sigma(\mathfrak u):=(\mathfrak u_A,\mathfrak u_B),
		      \]
		      is bijective.
	\end{enumerate}
\end{theorem}
\begin{proof}
	\emph{(a)\(\Leftrightarrow\)(b).}
	This is \cref{thm:boolean-complete}.
	\medskip
	\emph{(b)\(\Rightarrow\)(c).}
	Let \(\mathcal T\) be as in \cref{prop:tensor-exists}. By \emph{(F2)},
	every \(E\in\B\) is a finite union of intersections \(A\cap B\), so
	\[
		\B\subseteq\mathcal T.
	\]
	Conversely, every generator \(A\cap B\) of \(\mathcal T\) already
	belongs to \(\B\), so \(\mathcal T\subseteq\B\). Hence
	\[
		\B=\mathcal T.
	\]
	Therefore \(\B\) satisfies the universal property of \(\BA\otimes\BB\).
	\medskip
	\emph{(c)\(\Rightarrow\)(d).}
	Assume \(\B=\BA\otimes\BB\), via a fixed identification.
	First, \(\Sigma\) is well defined by \cref{lem:marginals}.
	\smallskip
	\emph{Injectivity.}
	Suppose
	\[
		\Sigma(\mathfrak u)=\Sigma(\mathfrak v).
	\]
	Since \(\B\) is generated by the rectangle generators
	\[
		A\otimes B,\qquad A\in\BA,\ B\in\BB,
	\]
	it suffices to show agreement on these generators. For any ultrafilter
	\(\mathfrak u\),
	\[
		A\otimes B\in\mathfrak u
		\iff
		A\in\mathfrak u_A \text{ and } B\in\mathfrak u_B.
	\]
	The right-hand side depends only on \(\Sigma(\mathfrak u)\). Thus
	\(\mathfrak u\) and \(\mathfrak v\) agree on all generators, hence on
	all of \(\B\). So \(\Sigma\) is injective.
	\smallskip
	\emph{Surjectivity.}
	Let
	\[
		(\mathfrak a,\mathfrak b)\in\UF(\BA)\times\UF(\BB).
	\]
	Consider the family
	\[
		\mathcal G(\mathfrak a,\mathfrak b)
		:=
		\{A\cap B:\ A\in\mathfrak a,\ B\in\mathfrak b\}\subseteq\B.
	\]
	We claim that \(\mathcal G(\mathfrak a,\mathfrak b)\) has the finite
	intersection property. Let
	\[
		A_1\cap B_1,\dots,A_n\cap B_n\in\mathcal G(\mathfrak a,\mathfrak b).
	\]
	Then
	\[
		A^\ast:=\bigcap_{i=1}^n A_i\in\mathfrak a,
		\qquad
		B^\ast:=\bigcap_{i=1}^n B_i\in\mathfrak b.
	\]
	Since ultrafilters do not contain \(0\), both \(A^\ast\) and \(B^\ast\)
	are nonempty. By \emph{(F1)},
	\[
		A^\ast\cap B^\ast\neq\varnothing.
	\]
	Hence
	\[
		\bigcap_{i=1}^n (A_i\cap B_i)=A^\ast\cap B^\ast\neq\varnothing.
	\]
	So \(\mathcal G(\mathfrak a,\mathfrak b)\) has the finite intersection
	property. By the ultrafilter extension theorem, it extends to some
	\(\mathfrak u\in\UF(\B)\). By construction,
	\[
		\mathfrak u_A=\mathfrak a,
		\qquad
		\mathfrak u_B=\mathfrak b.
	\]
	Thus \(\Sigma\) is surjective.
	\medskip
	\emph{(d)\(\Rightarrow\)(b).}
	Assume \(\Sigma\) is bijective.
	To prove \emph{(F1)}, let \(A\in\BA\) and \(B\in\BB\) be nonempty.
	Choose ultrafilters \(\mathfrak a\in\UF(\BA)\) and \(\mathfrak b\in\UF(\BB)\)
	such that
	\[
		A\in\mathfrak a,
		\qquad
		B\in\mathfrak b.
	\]
	By surjectivity of \(\Sigma\), there exists \(\mathfrak u\in\UF(\B)\)
	with
	\[
		\Sigma(\mathfrak u)=(\mathfrak a,\mathfrak b).
	\]
	Then \(A,B\in\mathfrak u\), hence
	\[
		A\cap B\in\mathfrak u.
	\]
	Since an ultrafilter cannot contain \(\varnothing\), it follows that
	\[
		A\cap B\neq\varnothing.
	\]
	Thus \emph{(F1)} holds.
	To prove \emph{(F2)}, let
	\[
		\B_{\mathrm{rect}}
		:=
		\Bool\bigl(\{A\cap B:\ A\in\BA,\ B\in\BB\}\bigr)
		\subseteq \B.
	\]
	Assume for contradiction that
	\[
		\B_{\mathrm{rect}}\neq \B.
	\]
	Then there exists \(E\in\B\setminus\B_{\mathrm{rect}}\). Since
	\(\B_{\mathrm{rect}}\) is a proper Boolean subalgebra of \(\B\), there
	exist distinct ultrafilters \(\mathfrak u_1,\mathfrak u_2\in\UF(\B)\)
	which agree on \(\B_{\mathrm{rect}}\) but disagree on \(E\). In
	particular they agree on \(\BA\) and on \(\BB\), since both
	subalgebras lie in \(\B_{\mathrm{rect}}\). Therefore
	\[
		\Sigma(\mathfrak u_1)=\Sigma(\mathfrak u_2),
	\]
	contradicting injectivity of \(\Sigma\). Hence
	\[
		\B=\B_{\mathrm{rect}},
	\]
	which is exactly \emph{(F2)}.
\end{proof}
\begin{remark}
	No appeal to Stone duality is required in the proof of
	\cref{thm:ultrafilter-splitting}. The only external inputs are the
	ultrafilter extension theorem and the compactness results of
	\cref{ch:foundation}.
\end{remark}
% ============================================================
% GLOBAL CLASSIFICATION
% ============================================================
\begin{theorem}[Rigidity of comparison completeness]
	\label{thm:comparison-completeness-rigidity}
	Under the standing principle \cref{sp:closed-world}, let
	\((U,\mathcal C)\) be a comparison world in the local distinguishability
	setup
	\[
		\alpha\cap\beta=\Delta_U.
	\]
	Then the following are equivalent:
	\begin{enumerate}[label=\textup{(\arabic*)}]
		\item \((U,\mathcal C)\) is rectangularly complete.
		\item The canonical factor map
		      \[
			      \Theta:U\to X_A\times X_B
		      \]
		      is bijective.
		\item There exists a product presentation
		      \[
			      U\cong Y_A\times Y_B
		      \]
		      whose coordinate congruences coincide with \(\alpha\) and \(\beta\).
		\item The comparison Boolean algebra satisfies the rectangular
		      factorization property.
		\item The comparison Boolean algebra decomposes as
		      \[
			      \B\cong\BA\otimes\BB.
		      \]
		\item The marginal ultrafilter map
		      \[
			      \Sigma:\UF(\B)\to\UF(\BA)\times\UF(\BB),
			      \qquad
			      \Sigma(\mathfrak u):=(\mathfrak u_A,\mathfrak u_B),
		      \]
		      is bijective.
	\end{enumerate}
	When these conditions hold, the product decomposition
	\[
		U\cong X_A\times X_B
	\]
	is canonical, every realizing product presentation factors uniquely
	through \(\Theta\), the symmetry group \(G=\Aut(U,\mathcal C)\) acts
	diagonally on \(X_A\times X_B\), and the quotient data
	\[
		\Phys=(X_A\times X_B)/G
	\]
	are determined.
\end{theorem}
\begin{proof}
	The equivalence of \textup{(1)}, \textup{(2)}, and \textup{(3)} is
	\cref{thm:classification}. The equivalence of \textup{(1)} and
	\textup{(4)} is \cref{thm:boolean-complete}. The equivalence of
	\textup{(1)}, \textup{(4)}, \textup{(5)}, and \textup{(6)} is
	\cref{thm:ultrafilter-splitting}. The uniqueness and rigidity
	statements follow from \cref{prop:terminality,thm:diagonal,cor:terminal-data}.
\end{proof}
% ============================================================
\section{Structural completion of comparison worlds}
\label{sec:part1-structural-completion}
Starting from a comparison world \((U,\mathcal C)\), this chapter has
identified the intrinsic congruences \(\alpha\) and \(\beta\), the
canonical quotient sets
\[
	X_A:=U/\alpha,
	\qquad
	X_B:=U/\beta,
\]
and the factor map
\[
	\Theta:U\to X_A\times X_B,
	\qquad
	\Theta(u)=([u]_\alpha,[u]_\beta).
\]
Building on the compactness boundary from \Cref{ch:foundation}, the
chapter has shown that local intrinsic distinguishability together with
rectangular comparison completeness rigidly determines the global
comparison structure.

\begin{theorem}[Structural completion of comparison worlds]
	\label{thm:part1-structural-completion}
	Under the standing principle \cref{sp:closed-world}, let
	$(U,\mathcal C)$ be a comparison world satisfying the compactness
	boundary condition of \cref{sp:compactness}, the local
	distinguishability condition
	\[
		\alpha\cap\beta=\Delta_U,
	\]
	and the closure condition of \cref{sp:closure}, equivalently
	rectangular comparison completeness.
	Then the following structure is canonically determined.
	\begin{enumerate}[label=\textup{(\roman*)}]
		\item
		      The canonical factor map
		      \[
			      \Theta:U\to X_A\times X_B
		      \]
		      is bijective. Hence
		      \[
			      U\cong X_A\times X_B .
		      \]
		\item
		      The intrinsic symmetry group
		      \[
			      G:=\Aut(U,\mathcal C)
		      \]
		      acts diagonally on
		      \[
			      X:=X_A\times X_B .
		      \]
		\item
		      The comparison world therefore determines the orbit projection
		      \[
			      \pi:X\to \Phys:=X/G .
		      \]
	\end{enumerate}
\end{theorem}

\begin{proof}
	Item \textup{(i)} is the comparison completeness classification theorem
	(\cref{thm:classification}).
	Item \textup{(ii)} is the diagonal action theorem
	(\cref{thm:diagonal}).
	Item \textup{(iii)} follows from the canonical orbit construction
	(\cref{cor:terminal-data}).
\end{proof}

The same closure condition has already been expressed in three
equivalent Chapter 2 languages by
\cref{thm:comparison-completeness-rigidity}:
comparison-theoretic, Boolean-algebraic, and ultrafilter-theoretic.
Thus the canonical product decomposition emerges simultaneously as
relational rigidity, Boolean factorization, and ultrafilter splitting,
without introducing any semantic content beyond the orbit quotient
already fixed by the comparison data.

\begin{remark}[Interpretation]
	The theorem shows that once Boolean compactness, intrinsic
	distinguishability, and rectangular comparison completeness hold, the
	global structure of the comparison world is rigid. The state space
	admits a canonical product decomposition, the intrinsic symmetry group
	acts diagonally on that product, and the orbit quotient \(\Phys\) is
	thereby fixed by the comparison data themselves.
	Quotient semantics is the next chapter's semantic reading of this
	determined quotient backbone, not an additional structure imposed from
	outside.
\end{remark}

\begin{remark}[Role in the manuscript]
	\cref{thm:part1-structural-completion} closes Part~I's comparison
	stage by fixing the canonical product, diagonal symmetry, and orbit
	quotient data, together with their equivalent Boolean and ultrafilter
	formulations already proved in this chapter.
	\Cref{ch:rotation} then starts from this quotient backbone to
	establish closed-system quotient semantics and the
	subsystem-attribution boundary.
\end{remark}

\section{Conclusion}
\label{sec:bottom-transition}
Part~I closes with a rigidity theorem of exact scope: rectangular
comparison completeness determines canonical product structure, diagonal
symmetry, the canonical orbit quotient, and the matching Boolean and
ultrafilter formulations without residual structural freedom.
\Cref{ch:rotation} starts from this fixed quotient backbone and
establishes quotient descent before classifying the remaining
admissible non-quotient enrichment loci.